\documentclass[11pt]{article}
\usepackage[margin=0.85in]{geometry}
\usepackage{lmodern}
\usepackage{microtype}
\usepackage{enumitem}
\usepackage{listings}
\usepackage[T1]{fontenc}

\lstset{
  basicstyle=\ttfamily\small,
  columns=fullflexible,
  keepspaces=true,
  showstringspaces=false,
  breaklines=true,
  frame=single,
  xleftmargin=0.4em,
  xrightmargin=0.4em,
  aboveskip=0.6em,
  belowskip=0.6em
}

\pagestyle{empty}
\setlength{\parindent}{0pt}
\setlist[enumerate]{leftmargin=1.4em, itemsep=0.7em, topsep=0.4em}
\setlist[itemize]{leftmargin=1.4em, itemsep=0.35em, topsep=0.25em}

\begin{document}

\begin{center}
{\LARGE\bfseries Sample for Quiz 1}\\[0.35em]
DATA 412/612
\end{center}

\vspace{0.5em}

\begin{enumerate}
\item Create the following matrix:
\begin{lstlisting}
A <- 10*t(matrix(99:0, 10,10))
A
\end{lstlisting}
\begin{lstlisting}
      [,1] [,2] [,3] [,4] [,5] [,6] [,7] [,8] [,9] [,10]
 [1,]  990  980  970  960  950  940  930  920  910   900
 [2,]  890  880  870  860  850  840  830  820  810   800
 [3,]  790  780  770  760  750  740  730  720  710   700
 [4,]  690  680  670  660  650  640  630  620  610   600
 [5,]  590  580  570  560  550  540  530  520  510   500
 [6,]  490  480  470  460  450  440  430  420  410   400
 [7,]  390  380  370  360  350  340  330  320  310   300
 [8,]  290  280  270  260  250  240  230  220  210   200
 [9,]  190  180  170  160  150  140  130  120  110   100
[10,]   90   80   70   60   50   40   30   20   10     0
\end{lstlisting}

\item Find the sum of the entries of the even columns of the above matrix.
\begin{lstlisting}
sum(A[,2],A[,4],A[,6],A[,8],A[,10])
# [1] 24500
\end{lstlisting}

\item Suppose you have the following data frame:
\begin{lstlisting}
df <- data.frame(
  Name = c("Johny", "Alisson", "Andy", "Sarah", "Alex"),
  Age = c(25, 30, 28, 22, 32),
  Annual_Income = c(50000, 60000, 55000, 48000, 70000),
  Account_Balance = c(1700, 2000, 1800, 900, 2500.90)
)
df
\end{lstlisting}
\begin{lstlisting}
     Name Age Annual_Income Account_Balance
1   Johny  25         50000          1700.0
2 Alisson  30         60000          2000.0
3    Andy  28         55000          1800.0
4   Sarah  22         48000           900.0
5    Alex  32         70000          2500.9
\end{lstlisting}
\begin{itemize}
\item What is the age average of people who get more than \$51k in a year.
\begin{lstlisting}
mean(df$Age[df$Annual_Income > 51000])
# [1] 30
\end{lstlisting}
\item What are their names?
\begin{lstlisting}
df$Name[df$Annual_Income > 51000]
# [1] "Alisson" "Andy"    "Alex"
\end{lstlisting}
\end{itemize}

\item Given a data frame \texttt{df} with only positive values, write a function that check if the number of columns in it is less than equal to 1 or it has negative values at all. If yes, it returns $-1$. If not, the function compares the geometric means (using the pipe operator) of the first and the second columns and produces 1 if that is correct, otherwise 0. Try your function on these inputs \texttt{data.frame(a=c(9,1,8))}, \texttt{data.frame(a=c(9,9,8), b=c(9,7,6))}, and \texttt{data.frame(a=c(9,1,8), b=c(9,17,6))}.
\begin{lstlisting}
library(tidyverse)

my_df_function <- function(df) {
  if (length(df) <= 1 | any(df < 0)) {
    return(-1)
  } else {
    df[,1] %>%
      log() %>%
      mean() %>%
      exp() -> geo_mean_1

    df[,2] %>%
      log() %>%
      mean() %>%
      exp() -> geo_mean_2

    if (geo_mean_1 >= geo_mean_2) {
      return(1)
    } else {
      return(0)
    }
  }
}

df <- data.frame(a=c(-9,1,8), b=c(-4,3,18))
my_df_function(df)   # [1] -1

df1 <- data.frame(a=c(9,9,8), b=c(9,7,6))
my_df_function(df1)  # [1] 1

df2 <- data.frame(a=c(9,1,8), b=c(9,17,6))
my_df_function(df2)  # [1] 0
\end{lstlisting}
\end{enumerate}

\end{document}
